\documentclass[aspectratio=169,11pt]{beamer}

\usepackage[T1]{fontenc}
\usepackage[utf8]{inputenc}
\usepackage[brazil]{babel}
\usepackage{booktabs}
\usepackage{xcolor}

\usetheme{default}
\usecolortheme{seahorse}
\setbeamertemplate{navigation symbols}{}
\setbeamertemplate{footline}[frame number]
\setbeamertemplate{itemize items}[circle]
\definecolor{ufla}{rgb}{0.09,0.30,0.20}
\setbeamercolor{structure}{fg=ufla}
\setbeamerfont{frametitle}{size=\large,series=\bfseries}

\title{Detecção de fraudes financeiras}
\subtitle{Agrupamento HDBSCAN em comparação com Random Forest e XGBoost}
\author{Heitor Rodrigues Sabino}
\institute{Universidade Federal de Lavras \\ \smallskip
\footnotesize Orientadora: Elaine Cecilia Gatto \\
Coorientador: Natanael Fabrício Dacioli Batista}
\date{2026}

\begin{document}

\begin{frame}
    \titlepage
\end{frame}

% =============================================================================
\begin{frame}{O problema}

    Fraude em cartão de crédito é um evento \textbf{raro} e \textbf{mutável}.

    \bigskip
    \begin{itemize}
        \item A base do Machine Learning Group da ULB tem 284.807 transações e
              apenas \textbf{492 fraudes} --- 0,172\%.
        \item Modelos supervisionados aprendem os padrões de fraude
              \textbf{já rotulados}. Fraudes novas, por definição, não estão nos
              rótulos.
        \item Rotular transações é caro e sempre chega atrasado.
    \end{itemize}

    \bigskip
    \begin{block}{Hipótese}
        Se fraude é uma anomalia, um algoritmo de densidade deveria encontrá-la
        \textbf{sem usar os rótulos} --- como ruído.
    \end{block}

\end{frame}

% =============================================================================
\begin{frame}{A pergunta de pesquisa}

    \begin{center}
    \large
    Um algoritmo não supervisionado de agrupamento por densidade compete com
    classificadores supervisionados na detecção de fraudes?
    \end{center}

    \bigskip
    \begin{columns}[t]
        \begin{column}{0.48\textwidth}
            \textbf{Não supervisionado}
            \begin{itemize}
                \item HDBSCAN
                \item Fraude = ruído (rótulo $-1$)
                \item Não vê a coluna \textit{Class} no treino
            \end{itemize}
        \end{column}
        \begin{column}{0.48\textwidth}
            \textbf{Supervisionados (referência)}
            \begin{itemize}
                \item Random Forest
                \item XGBoost
                \item Treinam com os rótulos
            \end{itemize}
        \end{column}
    \end{columns}

\end{frame}

% =============================================================================
\begin{frame}{Como foi resolvido --- validação}

    \textbf{Validação cruzada aninhada estratificada}

    \bigskip
    \begin{center}
    \begin{tabular}{lrr}
        \toprule
        \textbf{Partição} & \textbf{Treino} & \textbf{Avaliação} \\
        \midrule
        5 \textit{folds} externos (80/20) & 227.845 & 56.961 \\
        Divisão interna (80/20)           & 182.276 & 45.569 \\
        \bottomrule
    \end{tabular}
    \end{center}

    \bigskip
    \begin{itemize}
        \item Estratificada: cada partição preserva os 0,172\% de fraude.
        \item Os dados que escolhem hiperparâmetros \textbf{nunca} são os que
              estimam o desempenho.
    \end{itemize}

\end{frame}

% =============================================================================
\begin{frame}{Como foi resolvido --- pré-processamento e busca}

    \textbf{Pré-processamento}
    \begin{itemize}
        \item \texttt{Time}: padronização Z-score --- variável incremental e
              limitada.
        \item \texttt{Amount}: \textit{Robust Scaler} --- os valores vão de
              0,00 a 25.691,16, com mediana de 22,00.
        \item \texttt{V1}--\texttt{V28}: mantidas; já são componentes de PCA.
    \end{itemize}

    \bigskip
    \textbf{Busca em grade no ciclo interno}
    \begin{itemize}
        \item Supervisionados: seleção por \textbf{F1}.
        \item HDBSCAN: seleção por \textbf{DBCV}
              (\textit{Density-Based Clustering Validation}).
    \end{itemize}

    \bigskip
    \begin{alertblock}{Por que DBCV}
        É uma métrica interna: avalia a qualidade do agrupamento sem consultar
        os rótulos. Usar F1 aqui contaminaria a etapa não supervisionada.
    \end{alertblock}

\end{frame}

% =============================================================================
\begin{frame}{Resultado preliminar --- e o problema que ele revelou}

    O HDBSCAN é \textbf{instável entre partições quase idênticas}.

    \bigskip
    Mesma configuração (\texttt{min\_cluster\_size}=5,
    \texttt{min\_samples}=300), cinco \textit{folds} internos:

    \begin{center}
    \begin{tabular}{lrrr}
        \toprule
         & \textbf{mínimo} & \textbf{máximo} & \textbf{amplitude} \\
        \midrule
        DBCV                & 0,107 & 0,959 & 0,852 \\
        Pontos de ruído     & 7.985 & 130.231 & 122.246 \\
        \bottomrule
    \end{tabular}
    \end{center}

    \bigskip
    \begin{itemize}
        \item Como fraude é o rótulo de ruído, essa amplitude determina
              diretamente quantas transações o modelo acusa.
        \item O F1 do HDBSCAN ficou entre \textbf{0,011 e 0,070}.
    \end{itemize}

\end{frame}

% =============================================================================
\begin{frame}{Questão em aberto}

    \begin{block}{Critério de seleção: máximo ou média?}
        O ciclo interno elege o maior DBCV entre os cinco \textit{folds}. Numa
        distribuição tão instável, o máximo seleciona também o \textit{fold}
        mais favorável. A média é a escolha usual da validação aninhada.
    \end{block}

\end{frame}

\end{document}
